\documentclass[11pt]{article}
\usepackage[a4paper,margin=25mm]{geometry}
\usepackage{amsmath,amssymb,amsthm}
\usepackage[hidelinks]{hyperref}
\newtheorem{theorem}{Theorem}
\newtheorem{lemma}[theorem]{Lemma}
\theoremstyle{remark}
\newtheorem*{remark}{Remark}
\newcommand{\hypo}{\equiv}
\title{Hypoplactic Classes Are Not Counted by Plane Partitions:\\
A Disproof of Conjecture 00000003842}
\author{GitHub: gaochengzhecpu}
\date{}
\begin{document}
\maketitle
\begin{abstract}
The source asserts that the number of hypoplactic classes of words of
length at most $\ell$ over an $n$-letter alphabet equals the number of
plane partitions in a $2\times n\times\ell$ box. This is false. Over two
letters there are $7$ classes of length at most $2$ and $4$ of length
exactly $2$, but $20$ plane partitions in the $2\times2\times2$ box. The
identity also fails for every $\ell\ge1$ when $n=1$ and for every
$n\ge1$ when $\ell=1$. The hypoplactic monoid, its classes, plane
partitions, and the negation of the identity are formalised in Lean. An
earlier submission had the right example but no formal content; its
error is described below.
\end{abstract}

\section*{Definitions and reading of the source}
Let $A=\{1<2<\dots<n\}$ and let $A^*$ be the free monoid of words. The
\emph{plactic congruence} is generated by the Knuth relations
\[
 acb\equiv cab\ \ (a\le b<c),\qquad bac\equiv bca\ \ (a<b\le c).
\]
The \emph{hypoplactic congruence} $\hypo$ is generated by the Knuth
relations together with the quartic relations
\[
 cadb\hypo acbd\ \ (a\le b<c\le d),\qquad
 bdac\hypo dbca\ \ (a<b\le c<d).
\]
The hypoplactic monoid is $A^*/{\hypo}$; it is a quotient of the plactic
monoid (Krob--Thibon, Novelli). All defining relations preserve length,
so every class has a length.

A \emph{plane partition in an $a\times b\times c$ box} is an $a\times b$
array of integers in $\{0,\dots,c\}$ that is weakly decreasing along
every row and every column. Their number $\mathrm{PP}(a,b,c)$ is
symmetric in $a,b,c$, so the order of the three sides in the source does
not matter.

The English text counts the classes ``of length at most $\ell$''. The
Chinese text counts the classes ``of length $\ell$''. We treat both.
Write $h(n,\ell)$ for the number of classes of words of length exactly
$\ell$, and $H(n,\ell)=\sum_{m\le\ell}h(n,m)$ for the number of classes
of words of length at most $\ell$, the empty word included. The source
claims $H(n,\ell)=\mathrm{PP}(2,n,\ell)$, or in the Chinese wording
$h(n,\ell)=\mathrm{PP}(2,n,\ell)$, for all $n$ and $\ell$.

\section*{The disproof}
\begin{lemma}\label{lem:short}
If $u\hypo v$ then $u$ and $v$ have the same length. If moreover
$u\ne v$, then this length is at least $3$.
\end{lemma}
\begin{proof}
Let $R$ be the relation on $A^*$ defined by: $u\mathrel{R}v$ if $u=v$,
or $u$ and $v$ have the same length and that length is at least $3$.
It is an equivalence relation. It is compatible with products: if
$u\mathrel{R}v$ and $u'\mathrel{R}v'$, then either $u=v$ and $u'=v'$,
so $uu'=vv'$, or one of the two pairs consists of words of length at
least $3$, and then $uu'$ and $vv'$ have the same length, at least $3$.
So $R$ is a congruence. Both sides of every defining relation have
length $3$ or $4$, so $R$ contains the defining relations. Hence $R$
contains $\hypo$.
\end{proof}

\begin{theorem}\label{thm:main}
\begin{enumerate}
\item[(a)] For $\ell\le2$, $h(n,\ell)=n^\ell$. Hence $H(n,1)=n+1$ and
$H(n,2)=n^2+n+1$.
\item[(b)] $h(1,\ell)=1$ and $H(1,\ell)=\ell+1$ for every $\ell$.
\item[(c)] $\mathrm{PP}(2,1,\ell)\ge\ell+2$ for $\ell\ge1$,
$\mathrm{PP}(2,n,1)\ge n+2$ for $n\ge1$, and
$\mathrm{PP}(2,2,2)=20$.
\end{enumerate}

\noindent Consequently
\[
 h(2,2)=4,\quad H(2,2)=7,\quad \mathrm{PP}(2,2,2)=20;
\]
for every $\ell\ge1$, $h(1,\ell)<H(1,\ell)<\mathrm{PP}(2,1,\ell)$; and
for every $n\ge1$, $h(n,1)<H(n,1)<\mathrm{PP}(2,n,1)$. Both readings of
the conjectured identity are false.
\end{theorem}
\begin{proof}
(a) By Lemma~\ref{lem:short} a word of length at most $2$ is alone in
its class. So the classes of length $\ell\le2$ correspond bijectively to
the $n^\ell$ words of length $\ell$. Classes of different lengths are
different, which gives the values of $H$.

(b) Over one letter there is exactly one word of each length, and by
Lemma~\ref{lem:short} words of different lengths are inequivalent.

(c) In the $2\times1\times\ell$ box the $\ell+1$ constant arrays and
the array with entries $1$ above $0$ are $\ell+2$ different plane
partitions when $\ell\ge1$. In the $2\times n\times1$ box, take for
$k=0,\dots,n$ the array whose first row is $1^k0^{n-k}$ and whose second
row is zero, and the array with all entries $1$; these are $n+2$
different plane partitions when $n\ge1$. A plane partition in the
$2\times2\times2$ box is an array
$\bigl(\begin{smallmatrix}x&y\\z&w\end{smallmatrix}\bigr)$ with
$2\ge x\ge y\ge w\ge0$ and $x\ge z\ge w$. For fixed $x$ and $w$ there
are $(x-w+1)^2$ choices of $(y,z)$, so the number is
$1+(4+1)+(9+4+1)=20$.

The displayed consequences follow: $4<7<20$; $1<\ell+1<\ell+2$ for
$\ell\ge1$; and $n<n+1<n+2$.
\end{proof}

\begin{remark}
The exact values are $\mathrm{PP}(2,1,\ell)=\binom{\ell+2}2$ and
$\mathrm{PP}(2,n,1)=\binom{n+2}2$, and by MacMahon's formula
$\mathrm{PP}(2,2,\ell)=(\ell+1)(\ell+2)^2(\ell+3)/12$. If the empty word
is not counted, $H$ decreases by one and the inequalities above remain
strict. For orientation: hypoplactic classes correspond to quasi-ribbon
tableaux, which gives
$h(n,\ell)=\sum_{k\ge1}\binom{\ell-1}{k-1}\binom{n+\ell-k}{\ell}$ for
$\ell\ge1$, for example $h(2,\ell)=2\ell$. For $n=2$ this gives:
\[
\begin{array}{c|cccccc}
 \ell&0&1&2&3&4&5\\\hline
 h(2,\ell)&1&2&4&6&8&10\\
 H(2,\ell)&1&3&7&13&21&31\\
 \mathrm{PP}(2,2,\ell)&1&6&20&50&105&196
\end{array}
\]
So $H(2,\ell)$ grows like $\ell^2$ and $\mathrm{PP}(2,2,\ell)$ like
$\ell^4/12$; the two sides agree only in the trivial case $\ell=0$. This
remark is context. It is not formalised and the disproof does not depend
on it.
\end{remark}

\section*{The earlier submission and its error}
Pull request 280 of the competition repository (by
\texttt{orionsheep}) proposed the instance $n=2$, $\ell=2$ with the
values $6$ and $20$. It was closed without merging on 3 October 2026.
The reviewer's stated reason was that the main Lean theorem was
arithmetic with none of the objects of the conjecture: it asserted
$2+4=6$, $6\ge6$, $17280/864=20$ and $20>6$; no hypoplactic class and no
plane partition was defined in Lean; and the word count was entered by
hand and omitted the empty word (there are $7$ words of length at most
$2$, not $6$).

That is the error. Its final Lean theorem \texttt{conjecture\_refuted}
is the conjunction of these four numerical facts. The statements
``there are $6$ words'', ``the classes are at most as many as the
words'' and ``the box contains $20$ plane partitions'' appear only in
comments; the formal part says nothing about words, congruence classes
or plane partitions, so it does not certify the counterexample. The
mathematical idea was right: the classes of nonempty words of length at
most $2$ over two letters are at most $6$ in number, the box contains
$20$ plane partitions, and so the instance $n=2$, $\ell=2$ refutes the
identity, as the reviewer also noted.

The present submission defines the hypoplactic congruence on Mathlib's
free monoid, the sets of classes of given length, and plane partitions
in a box; it proves the exact counts $4$, $7$ and $20$, and negates the
identity as a statement quantified over all $n$ and $\ell$.

\section*{Formal verification and scope}
\texttt{KnuthRel} and \texttt{QuarticRel} are the four families of
relations above on Mathlib's \texttt{FreeMonoid} of a linearly ordered
alphabet. \texttt{hypoCon} is the congruence generated by them
(Mathlib's \texttt{conGen}), and \texttt{Hypo} is the quotient monoid.
\texttt{placticCon\_le\_hypoCon} states that Knuth equivalence is
contained in it. The theorems \texttt{knuth\_instance} and
\texttt{quartic\_instance} show $010\hypo100$ and $1010\hypo0101$ over
two letters, so the relations are really present. Lemma~\ref{lem:short}
is
\begin{center}\small
\texttt{length\_eq\_of\_hypoCon},\quad
\texttt{eq\_or\_three\_le\_of\_hypoCon}.
\end{center}

The sets of elements of the hypoplactic monoid represented by a word of
length exactly $\ell$, respectively at most $\ell$, are
\begin{center}\small
\texttt{classesOfLength},\quad \texttt{classesUpToLength};
\end{center}
their sizes are $h$ and $H$. Part (a) is
\begin{center}\small
\texttt{ncard\_classesOfLength\_of\_le\_two},\quad
\texttt{ncard\_classesUpToLength\_two},
\end{center}
and part (b) is
\begin{center}\small
\texttt{ncard\_classesOfLength\_of\_unique},\quad
\texttt{ncard\_classesUpToLength\_of\_unique}.
\end{center}

\texttt{PlanePartition a b c} is the type of arrays indexed by
\texttt{Fin a} and \texttt{Fin b} with values in \texttt{Fin (c + 1)}
that are weakly decreasing along rows and columns. The count $20$ is
\texttt{card\_pp\_222}, a kernel evaluation over all $81$ arrays. The
two lower bounds of (c) are
\begin{center}\small
\texttt{card\_pp\_one\_column},\quad
\texttt{card\_pp\_height\_one}.
\end{center}

The proposition \texttt{ClaimedIdentityUpTo} is the identity
$H(n,\ell)=\mathrm{PP}(2,n,\ell)$ for all $n$ and $\ell$, and
\texttt{ClaimedIdentityExact} is the same with $h$. Their negations are
\begin{center}\small
\texttt{conjecture\_false\_upTo},\quad
\texttt{conjecture\_false\_exact}.
\end{center}
The theorems \texttt{two\_letters\_length\_two}, \texttt{one\_letter}
and \texttt{length\_one} are the three displayed consequences of
Theorem~\ref{thm:main}. No \texttt{sorry}, \texttt{native\_decide} or
custom axiom occurs.

Not formalised: the exact binomial values and MacMahon's formula, the
symmetry of $\mathrm{PP}$ in its three arguments, the quasi-ribbon
formula and the table in the remark, and the correspondence of the
presentation used here with the definition of the hypoplactic monoid
through quasi-ribbon tableaux. The counterexamples do not depend on the
exact form of the relations: they use only that every defining relation
preserves length and has length at least $3$.

\section*{Source and reproduction}
The exact bilingual source is included byte-for-byte as
\texttt{SOURCE.md}. Its repository path is
\begin{center}\small\texttt{conjectures/00000003842.md}.\end{center}
The project pins Lean 4.19.0 and Mathlib commit
\begin{center}\small\texttt{c44e0c8ee63ca166450922a373c7409c5d26b00b}.\end{center}
From \texttt{lean/}, run \texttt{lake build}, then
\begin{center}\small\texttt{lake env lean -DwarningAsError=true Main.lean}.\end{center}
The included public-Git manifest locks all dependency revisions. The
\texttt{verification/} directory records the fresh-build logs,
theorem-axiom reports, artifact hashes, review and final PDF
inspection. No supplementary numerical computation is required.
\end{document}
